\documentclass[11pt]{article}

\usepackage[utf8]{inputenc}
\usepackage[T1]{fontenc}
\usepackage{geometry}
\usepackage{amsmath}
\usepackage{amssymb}
\usepackage{booktabs}
\usepackage{hyperref}
\usepackage{xurl}
\usepackage{xcolor}
\usepackage{listings}
\usepackage{enumitem}
\usepackage{parskip}
\usepackage{graphicx}
\graphicspath{{../images/}}

\geometry{margin=1in}

\lstset{
  language=Java,
  basicstyle=\ttfamily\small,
  keywordstyle=\color{blue},
  commentstyle=\color{gray},
  stringstyle=\color{teal},
  showstringspaces=false,
  columns=fullflexible,
  frame=single,
  framerule=0.2pt,
  breaklines=true
}

\setlist[itemize]{topsep=0.3em,itemsep=0.2em}
\setlist[enumerate]{topsep=0.3em,itemsep=0.2em}

% Simple per-section source line.
\newcommand{\notesources}[1]{\par\noindent{\small\color{gray}\textit{Sources:} \sloppy #1\par}}

% Unit-wise source strings for this subject (used by slides/notes).
% Path is relative to the lecture `latex/` folder (the compilation CWD).
% OOPS with Java (CSEG1044) — unit-wise sources
%
% These are short, student-facing reference strings intended to be used with:
%   - \slidesources{...}  (in beamer slides)
%   - \notesources{...}   (in lecture notes)
%
% Style inspiration: other subjects in My-Subjects/ use semicolon-separated sources
% and include an "accessed" date for web documentation.

\newcommand{\OOPSUnitISources}{Schildt, \textit{Java: A Beginner's Guide} (10e), McGraw-Hill (Java fundamentals + OOP basics).; Oracle Java Tutorials: Getting Started + Language Basics + Classes and Objects (accessed \today).; Java SE API docs (e.g., \texttt{String}, \texttt{Scanner}) (accessed \today).}

\newcommand{\OOPSUnitIISources}{Schildt, \textit{Java: The Complete Reference} (12e), McGraw-Hill (classes, inheritance, packages, interfaces).; Bloch, \textit{Effective Java} (3e), Addison-Wesley, 2018 (design choices; interfaces vs abstract classes).; Oracle Java Tutorials: Classes and Objects + Interfaces + Packages (accessed \today).; Java SE API docs (e.g., \texttt{Comparable}, \texttt{Runnable}) (accessed \today).}

\newcommand{\OOPSUnitIIISources}{Schildt, \textit{Java: The Complete Reference} (12e) (nested/inner classes, exceptions, threads, I/O).; Oracle Java Tutorials: Nested Classes + Exceptions + Concurrency + Basic I/O (accessed \today).; Java SE API docs (e.g., \texttt{Thread}, \texttt{AutoCloseable}) (accessed \today).}

\newcommand{\OOPSUnitIVSources}{Schildt, \textit{Java: The Complete Reference} (12e) (generics, lambdas, Swing, JDBC).; Bloch, \textit{Effective Java} (3e) (generics + lambdas).; Oracle Java Tutorials: Generics + Lambda Expressions + Swing + JDBC Basics (accessed \today).; Java SE API docs (e.g., \texttt{java.util.function}, \texttt{javax.swing}, \texttt{java.sql}) (accessed \today).}

\newcommand{\OOPSUnitVSources}{Schildt, \textit{Java: The Complete Reference} (12e) (Collections Framework + wrapper classes).; Oracle Java docs: Collections Framework overview (accessed \today).; Java SE API docs (\texttt{java.util} collections; wrapper classes in \texttt{java.lang}) (accessed \today).}

\newcommand{\OOPSUnitVISources}{Git SCM: \textit{Pro Git} (2e) (version control workflow), accessed \today.; GitHub Docs (repositories, issues, pull requests), accessed \today.; Oracle Java Tutorials: Swing + JDBC (accessed \today).; JUnit 5 User Guide (accessed \today).; Apache Maven Project: Getting Started (accessed \today).; Gradle User Manual: Getting Started (accessed \today).; UML class diagrams: UML 2.x overview/spec (accessed \today).}

\newcommand{\OOPSMiscWhyInterfacesSources}{Schildt, \textit{Java: The Complete Reference} (12e) (interfaces).; Bloch, \textit{Effective Java} (3e) (prefer interfaces where appropriate).; Oracle Java Tutorials: Interfaces (accessed \today).; Java SE API docs (e.g., \texttt{Comparable}, \texttt{Runnable}) (accessed \today).}



\title{Object Oriented Programming (CSEG1044)\\Unit 01 -- Lecture 02 Notes\\Java Overview + Program Structure + Features}
\author{Tofik Ali}
\date{\today}

\begin{document}
\maketitle
\tableofcontents

\section{Lecture Overview}
This lecture makes you comfortable with the Java ``toolchain'' and the minimum you need to compile and run Java programs.
By the end, you should clearly understand:
\begin{itemize}
  \item what the \textbf{JDK}, \textbf{JRE}, and \textbf{JVM} mean,
  \item how \textbf{bytecode} enables Java to run on different platforms,
  \item the difference between \textbf{compile-time errors} and \textbf{runtime errors}.
\end{itemize}
\notesources{\OOPSUnitISources}

\section{Syllabus Alignment (Unit 01)}
\textbf{Unit 01:} Introduction to OOPs (Java overview + features)

\section{Learning Outcomes}
After this lecture, you should be able to:
\begin{itemize}
  \item Explain the Java ecosystem: JDK vs JRE vs JVM, and what bytecode is.
  \item Compile and run a Java program using \texttt{javac} and \texttt{java}.
  \item Identify whether an error is compile-time (during \texttt{javac}) or runtime (during \texttt{java}).
  \item List core features of Java and what they imply in practice.
\end{itemize}

\section{Key Terms (Quick)}
\begin{itemize}
  \item \textbf{JDK (Java Development Kit)}: tools for development (compiler \texttt{javac}, docs, etc.).
  \item \textbf{JRE (Java Runtime Environment)}: runtime needed to run Java apps (JVM + libraries).
  \item \textbf{JVM (Java Virtual Machine)}: executes Java bytecode (\texttt{.class}).
  \item \textbf{Bytecode}: platform-neutral instructions produced by \texttt{javac}.
  \item \textbf{Classpath}: where the JVM looks for compiled classes and libraries.
  \item \textbf{Compile-time error}: detected by the compiler before running.
  \item \textbf{Runtime error/exception}: occurs while the program is executing.
\end{itemize}

\section{Detailed Notes}
\subsection{Java ecosystem: JDK vs JRE vs JVM}
Think of Java in layers:
\begin{itemize}
  \item \textbf{JDK}: used by developers to \emph{create} programs. Includes \texttt{javac} (compiler) and other tools.
  \item \textbf{JRE}: used to \emph{run} programs. Includes the JVM and standard libraries.
  \item \textbf{JVM}: the engine that loads classes, verifies bytecode, and executes it.
\end{itemize}

\paragraph{Rule of thumb.} If you want to \emph{write/compile} Java, install the \textbf{JDK}. If you only want to \emph{run} Java (rare for students), the \textbf{JRE} is enough.

\subsection{Bytecode and the compile/run flow}
Java is often described as ``Write Once, Run Anywhere'' because compilation produces \textbf{bytecode}, not machine code.
The same \texttt{.class} file can run on Windows/Linux/macOS \emph{as long as there is a JVM for that OS}.

\textbf{Flow (mental model):}
\begin{enumerate}
  \item Write: \texttt{Hello.java}
  \item Compile: \texttt{javac Hello.java} $\rightarrow$ \texttt{Hello.class}
  \item Run: \texttt{java Hello} (JVM loads \texttt{Hello.class} and starts \texttt{main})
\end{enumerate}

\subsection{Basic program structure: class + \texttt{main}}
A minimal program has:
\begin{itemize}
  \item a \textbf{class} (container for code),
  \item the \textbf{entry point} method: \texttt{public static void main(String[] args)}.
\end{itemize}

\textbf{Important rule:} if you declare a \texttt{public} class, the file name must match it exactly.
Example: \texttt{public class Hello} must be in \texttt{Hello.java}.

\subsection{Compile-time vs runtime errors}
\textbf{Compile-time error} examples:
\begin{itemize}
  \item syntax errors (missing semicolon, wrong braces),
  \item type errors (assigning \texttt{String} to \texttt{int} without conversion),
  \item missing symbols (calling a method that doesn't exist).
\end{itemize}

\textbf{Runtime error/exception} examples:
\begin{itemize}
  \item \texttt{ArithmeticException} (divide by zero),
  \item \texttt{NullPointerException} (calling method on \texttt{null}),
  \item \texttt{ArrayIndexOutOfBoundsException} (invalid index).
\end{itemize}

\subsection{Features of Java (what they mean in practice)}
\begin{itemize}
  \item \textbf{Platform independent}: bytecode runs on any platform with a JVM.
  \item \textbf{Object-oriented (mostly)}: you model problems using classes/objects (but primitives exist for efficiency).
  \item \textbf{Automatic memory management}: garbage collector reclaims unused objects (no manual \texttt{free}).
  \item \textbf{Strong typing}: many errors are caught at compile time (good for large codebases).
  \item \textbf{Rich standard library}: collections, I/O, networking, GUI, database connectivity, etc.
  \item \textbf{Security model}: class loading + bytecode verification help reduce unsafe behavior.
\end{itemize}
\notesources{\OOPSUnitISources}

\section{Worked Example: Compile \& Run + Observe Errors}
\subsection*{A minimal program}
\begin{lstlisting}
public class Hello {
    public static void main(String[] args) {
        System.out.println("Hello, Java!");
    }
}
\end{lstlisting}

\subsection*{How to run (terminal)}
\begin{itemize}
  \item Compile: \texttt{javac Hello.java}
  \item Run: \texttt{java Hello}
\end{itemize}

\subsection*{Runtime exception example}
\begin{lstlisting}
public class DivideDemo {
    public static void main(String[] args) {
        int a = 10;
        int b = 0;
        System.out.println(a / b); // ArithmeticException at runtime
    }
}
\end{lstlisting}

\subsection*{Compile-time error example (explain in class)}
If you write \texttt{System.out.println("Hi")} without a semicolon, \texttt{javac} will stop and show an error.
This is a compile-time error: the program never runs.

\section{In-Class Exercises (with brief solutions)}
\begin{enumerate}
  \item \textbf{JDK/JRE/JVM:} Which one contains \texttt{javac}? \textbf{Answer:} JDK.
  \item \textbf{File naming:} What should the file name be if you write \texttt{public class Student}? \textbf{Answer:} \texttt{Student.java}.
  \item \textbf{Error type:} Is ``cannot find symbol'' a compile-time or runtime error? \textbf{Answer:} compile-time.
\end{enumerate}

\section{Self-Study Practice (no solutions)}
\begin{enumerate}
  \item Write \texttt{ArgsDemo} that prints all command line arguments with index numbers.
  \item Create 2 small programs: one that fails at compile time, one that fails at runtime. Write a 2-line explanation for each.
  \item Install the JDK and run \texttt{java -version} and \texttt{javac -version}. Note the output.
\end{enumerate}

\section{Checkpoint Questions (with sample answers)}
\textbf{Q1.} Why does Java use bytecode instead of compiling directly to machine code?\par
\textbf{Sample answer:} Bytecode makes the compiled program portable. The same \texttt{.class} can run on different operating systems because the JVM translates/executes bytecode for the current platform.

\vspace{0.6em}
\textbf{Q2.} What is the difference between a compile-time error and a runtime exception?\par
\textbf{Sample answer:} Compile-time errors are detected by \texttt{javac} before execution; runtime exceptions occur while the program is running after compilation succeeds.

\end{document}
